\section{Finite presentations and matching bounds}
\begin{lemma}\label{bound}
Let $W$ be a presented group, let $L\leq W$ have finite index $d$, and suppose that a presented group $T$ maps onto $L$. Suppose $T$ has a cyclic subgroup $C=\langle x\rangle$ of index $e$. If a complete integer relation matrix for $C/[C,C]$, on the specified generator $x$, has one column with nonzero gcd $m$, then
\[
 |W|\leq dem.
\]
If $W$ maps onto a finite marked group of order $dem$, that epimorphism is an isomorphism.
\end{lemma}
\begin{proof}
A group generated by one element is abelian. Thus $C=C/[C,C]\cong\mathbb Z/m\mathbb Z$: its defining relation subgroup in $\mathbb Z$ is generated by the matrix entries. Hence $|T|=em$, while surjectivity gives $|L|\leq em$. Coset partitioning now gives $|W|=d|L|\leq dem$. A surjection onto a group of this upper-bound order has trivial kernel.
\end{proof}
The hypothesis concerns a complete subgroup relation matrix, not a matrix from an arbitrary finite quotient. In our calculation it is computed by GAP's modified Todd--Coxeter algorithm for the subgroup generated by the actual element $x$ of the presented group. The simpler finite permutation image is used only on the opposite side of the inequality. GAP's exact presentation algorithms and ACE's completed coset enumeration are retained computational premises; neither a list of permutations satisfying the relators nor a closed table with no verified subgroup enumeration would replace them.

\begin{lemma}\label{prime}
Let $K$ be finite and let $B\subseteq K$. Suppose that for every prime $p\mid |K|$ and every conjugacy class of subgroups of order $p$, some member meets $B$. Then every nontrivial subgroup $N\leq K$ has a conjugate meeting $B$.
\end{lemma}
\begin{proof}
Cauchy's theorem gives a subgroup $C\leq N$ of prime order. By hypothesis $C^g\cap B\ne\varnothing$ for some $g$, and this intersection is contained in $N^g\cap B$.
\end{proof}
There is a complete finite test without enumerating every subgroup: form a cyclic subgroup from each element-class representative of order $p$, identify these subgroups up to conjugacy, and test their conjugates against the elements of $B$ of order $p$. Every subgroup of order $p$ occurs because any of its nonidentity elements generates it. This proof does not require all elements of order $p$ to be conjugate.

